\section{Evaluation Protocol}
\label{sec:evaluation}

All methods in this paper share the same \textbf{encoder architecture},
\textbf{view pipeline}, and \textbf{downstream metrics} so that differences
reflect the loss function rather than capacity or preprocessing.

\subsection{Architecture and training budget}
\label{sec:eval-arch}

\begin{itemize}
  \item \textbf{Encoder:} SmallResNet~\cite{he2016resnet} adapted to
        $32{\times}32$ inputs: three residual stages, global average pooling,
        linear projection to $d{=}256$.
  \item \textbf{Predictor:} two-layer MLP ($256 \rightarrow 512 \rightarrow 256$)
        with LayerNorm and GELU, used whenever the method aligns
        $g_\phi(f_\theta(x_c))$ to $f_\theta(\tilde{x})$.
  \item \textbf{Parameters:} $\sim$2.8M trainable (encoder + predictor);
        we refer to this as the \emph{3M} setting.
  \item \textbf{Optimization:} AdamW, learning rate $3{\times}10^{-4}$,
        weight decay $10^{-4}$, cosine decay, batch size $256$, seed $42$.
  \item \textbf{Primary budget:} \textbf{200 epochs} for the ablation in
        Section~\ref{sec:ablation}; additional runs at 100 epochs are
        reported where 200-epoch checkpoints are not yet available.
\end{itemize}

Unless noted, we use a \textbf{single online encoder} without an EMA teacher
(\texttt{use\_ema\_target=false}), matching the latent training modes in our
codebase.
EMA-JEPA results are reported separately as an upper reference with a
stop-gradient target encoder.

\subsection{View construction}
\label{sec:eval-views}

For every method:
\begin{enumerate}
  \item \textbf{Clean} $\tilde{x}$: random crop + horizontal flip (dataset
        default augmentation).
  \item \textbf{Corrupt} $x_c$: multi-block masking with ratio $0.6$ on a
        $4{\times}4$ patch grid (I-JEPA style).
  \item \textbf{Augmented} $x_a$ (when used): color jitter with
        brightness/contrast/saturation $=0.4$, hue $=0.1$.
\end{enumerate}
Methods that only use two views (e.g.\ pure VICReg on corrupt/clean) omit
$x_a$; dual-view InfoNCE and our hybrids use all three.

\subsection{Downstream evaluation}
\label{sec:eval-metrics}

We evaluate \textbf{only the frozen encoder} $f_\theta$; the predictor is
discarded at test time.

\paragraph{k-NN classification.}
Extract L2-normalized features from the training and test sets.
Classify each test image by majority vote among $k{=}20$ nearest training
neighbors under cosine distance.
k-NN stresses \emph{local} geometry and instance-level neighborhood structure.

\paragraph{Linear probe.}
Train a linear classifier (softmax cross-entropy, SGD, 50 epochs) on
\emph{unnormalized} train features; report top-1 accuracy on the test set.
The linear probe stresses \emph{global} class separability.

\paragraph{Reporting.}
We report the best test k-NN and the linear accuracy at the same checkpoint
unless stated otherwise.
For long runs we also plot both metrics every 10 epochs to visualize
convergence and the k-NN / linear trade-off
(Figure~\ref{fig:tradeoff}).

\subsection{Implementation}
\label{sec:eval-impl}

Training and evaluation are implemented in the public JEPA repository.
Each method corresponds to a JSON config under \texttt{configs/} with fixed
\texttt{training\_mode}, loss weights, and \texttt{run\_name}.
Checkpoints store model weights, optimizer state, and per-epoch metrics in
\texttt{history} for reproducible plotting.
